%-----------------------------QUESTÃO------------------------------------
%Cód.: TF3p24q07


\question [\questaopt] Os pontos de uma determinada região do espaço estão sob a
influência única de uma carga positiva pontual $ Q $. Sabe-se que em um
ponto $ A $, distante \SI{2}{\meter} da carga $ Q $, a intensidade do campo elétrico é
igual a $\SI[per-mode = symbol-or-fraction]{1,8e4}{\newton\per\coulomb}$. Determine:

\begin{parts}

\part o valor da carga elétrica $ Q $;

\part a intensidade do campo elétrico num ponto $ B $, situado a 			$\SI{30}{\centi \meter}$ da carga fonte $ Q $.\\
\textbf{Dados:} Constante eletrostática do vácuo $ \cteK $. 

\end{parts}


%--------------------------SOLUÇÃO---------------------------------------

\begin{EnvUplevel}
	\begin{solution}[2cm]
	
		(a)\\
		$$
		Q=\SI{+8,0}{\micro \coulomb}
		$$
			
		(b)
		$$
		E_B=\SI[per-mode=symbol-or-fraction]{8e5}{\newton \per \coulomb}
		$$

		
	\end{solution}
\end{EnvUplevel}
